%%%%%%%%%%  scan idx 68  |  folio 55  %%%%%%%%%%
vérifier que pour toute partie fermée irréductible $Y$ de $X\times k$, $\gamma(Y)$ est dans l'image de $p^!$. Il suffit de continuer le raisonnement en supposant $X$ irréductible~; si $p(Y)$ n'est pas dense dans $X$, notre assertion résulte de l'hypothèse de récurrence~; sinon, on a $Y=X\times k$ et l'assertion est triviale, ou bien $Y$ est une hypersurface dans $X\times k$. Soit alors $U$ un ouvert affine de $X$ formé de points non singuliers~; si $D$ est le diviseur sur $U\times k$ induit par $Y$, on sait (algèbre commutative) qu'il est linéairement équivalent à un diviseur de la forme $D'\times k$, où $D'$ est un diviseur sur $U$. Alors $\gamma_{U\times k}(D)-\gamma_{U\times k}(D'\times k)$ appartient à $\underline K(U\times k)^2$ (corollaire de la proposition 2.10), et d'après ce qu'on a déjà démontré, cet élément est de la forme $p^!(x')$ avec $x'\in\underline K(U)$. Donc $\gamma_{U\times k}(D)=p^!(y')$ avec $y'=x'+\gamma_U(D')$~; mais $y'$ est restriction d'un élément $y$ de $\underline K(X)$ (prop.\ 2.10), et par suite la restriction de $\gamma_{X\times k}(D)-p^!(y)$ à $U\times k$ est nulle~; d'après la proposition 2.10, $\gamma_{X\times k}(D)-p^!(y)$ provient d'un élément de $\underline K(X\times k-U\times k)=\underline K(X'\times k)$ avec $X'=X-U$. La démonstration s'achève en appliquant de nouveau l'hypothèse de récurrence.
\begin{flushright}cqfd.\end{flushright}

\textbf{Corollaire~:} \emph{Soient $X$ un espace algébrique et $p$ la projection de $X\times k^n$ sur $X$~; alors l'homomorphisme $p^!:\underline K(X)\longrightarrow\underline K(X\times k^n)$ est surjectif~; il est même bijectif si $X$ est non singulière.}

Récurrence sur $n$ pour la première assertion. Pour la deuxième, on pose $f(x)=(x,0)$ d'où $pf=\underline{1}$ et $f^!p^!=1$, ce qui prouve que $p^!$ est injectif.

\textbf{Proposition 2.10.} \emph{Soient $X$ un espace algébrique, $U$ une partie ouverte de $X$ et $Y$ son complémentaire dans $X$. Alors la suite d'homomorphismes naturels~:}
